\section{Isotope Half-Lives and a Common Survival Experiment}
\label{manual2:isotope-survival}

\subsection{Source records and nucleon bookkeeping}
\paragraph{Setup.}
The element reference remains the 118 placements and 126 slots of \secref{sec:coord0-enclosures}.
A separate isotope layer joins by atomic number and symbol. It retains the nuclear-state
assignment, selected evaluation, original notation, asymmetric uncertainty, decay channels
and experimental overlap. Total half-life, channel-partial half-life and an individual
waiting time are distinct quantities. Observational stability supplies no measured infinity.

\paragraph{Derivation.}
For the formal grouping of a record with proton count \(Z\), neutron count \(N=A-Z\), define
\begin{equation}
q=\min(\lfloor Z/2\rfloor,\lfloor N/2\rfloor),\quad
r_p=Z-2q,\quad r_n=N-2q,\qquad A=4q+r_p+r_n.
\label{eq:manual2-isotope-grouping}
\end{equation}
These are exact integer coordinates. Identifying a group with a physical alpha cluster
requires an additional construction. The bookkeeping alone supplies neither a loss
probability nor a half-life.

\paragraph{Specialization.}
Tables~\ref{tab:isotope-anchors}--\ref{tab:isotope-superheavy} keep the printed reference
compact. The companion retains the broader family records, source-specific updates and
uncertainty qualifications. H-1 and Cs-133 retain their recurrence roles. Fe-56 is the
selected binding benchmark; the Ni-62 digital control prevents calling Fe-56 the maximum
binding energy per nucleon. Pb-208 remains observationally stable, with its experimental
alpha-channel limit stored separately from total lifetime.

\begin{table}[htbp]
\centering\small
\begin{tabular}{lrrrl}
\toprule
Isotope & $Z,N$ & $q$ & $r_p,r_n$ & Adopted $T_{1/2}$ \\
\midrule
$^{1}\mathrm{H}$ & $1,0$ & $0$ & $1,0$ & observationally stable \\
$^{133}\mathrm{Cs}$ & $55,78$ & $27$ & $1,24$ & observationally stable \\
\bottomrule
\end{tabular}
\caption{Recurrence anchors. Stability is an observational classification; no infinite half-life is assigned.}
\label{tab:isotope-anchors}
\end{table}

\begin{table}[htbp]
\centering\small
\begin{tabular}{lrrrl}
\toprule
Isotope & $Z,N$ & $q$ & $r_p,r_n$ & Adopted $T_{1/2}$ \\
\midrule
$^{56}\mathrm{Fe}$ & $26,30$ & $13$ & $0,4$ & observationally stable \\
$^{208}\mathrm{Pb}$ & $82,126$ & $41$ & $0,44$ & observationally stable \\
\bottomrule
\end{tabular}
\caption{Binding and observational-stability controls. The digital reference retains Ni-62 and the channel-specific Pb-208 bound.}
\label{tab:isotope-benchmarks}
\end{table}

\begin{table}[htbp]
\centering\small
\begin{tabular}{lrrrl}
\toprule
Isotope & $Z,N$ & $q$ & $r_p,r_n$ & Adopted $T_{1/2}$ \\
\midrule
$^{294}\mathrm{Og}$ & $118,176$ & $59$ & $0,58$ & $0.58^{+0.44}_{-0.18}$ ms \\
$^{290}\mathrm{Lv}$ & $116,174$ & $58$ & $0,58$ & $8.3^{+3.5}_{-1.9}$ ms \\
$^{286}\mathrm{Fl}$ & $114,172$ & $57$ & $0,58$ & $166^{+40}_{-27}$ ms \\
$^{293}\mathrm{Lv}$ & $116,177$ & $58$ & $0,61$ & $95^{+63}_{-27}$ ms \\
$^{289}\mathrm{Fl}$ & $114,175$ & $57$ & $0,61$ & $2.4^{+0.8}_{-0.5}$ s \\
$^{298}\mathrm{Fl}$ & $114,184$ & $57$ & $0,70$ & unobserved target \\
\bottomrule
\end{tabular}
\caption{Selected superheavy records. Asymmetric uncertainties are expanded from the adopted last-digit notation. Fl-298 is an unobserved target. Newer experimental analyses remain separate in the companion.}
\label{tab:isotope-superheavy}
\end{table}

\begin{figure}[htbp]
\centering
\( {}^{294}\mathrm{Og}\xrightarrow{\alpha}{}^{290}\mathrm{Lv}
\xrightarrow{\alpha}{}^{286}\mathrm{Fl}\xrightarrow{\alpha}{}^{282}\mathrm{Cn}
\xrightarrow{\mathrm{SF}}\text{fragments} \)
\par\smallskip
\( {}^{293}\mathrm{Lv}\xrightarrow{\alpha}{}^{289}\mathrm{Fl}
\xrightarrow{\alpha}{}^{285}\mathrm{Cn} \)
\caption{Observed alpha relationships in the selected records. The first route is one branch: Fl-286 also fissions. Cn-285 has further alpha decay; it is a scope endpoint here, not an asserted terminal stable state.}
\label{fig:isotope-chains}
\end{figure}


\paragraph{Calibration.}
Evaluated half-lives retain their source units and explicit conversion to seconds.
Newer experimental analyses remain distinct from the adopted evaluation; overlapping
events are not averaged as independent data. Fl-298 has no measured half-life and receives
no empirical score. The progression Fl-289, Og-294, Fl-298 is a comparison set, not a decay
chain. The sourced alpha edges in Fig.~\ref{fig:isotope-chains} obey \(\Delta Z=-2\),
\(\Delta A=-4\); spontaneous fission does not acquire an invented unique daughter.

\paragraph{Falsification.}
Source, state, quantity and status mutations must fail without rewriting inputs.
The logarithmic half-life display names the selected isotope and gives limits, stability,
unknown records and uncovered elements nonnumeric treatments. It does not reinterpret
the structural-role colors or an arbitrary \(A\) arithmetic input as measured stability.

\aodprovenance{\path{manual2/data/m1/isotopes/} binds the source bytes, selected records,
updates and generated tables. The original element ledger and historical numeric records
remain unchanged.}
\aodliteraturenote{Evaluated isotope data: ENSDF~\cite{isotopes-ensdf}; evaluation cross-check:
NUBASE2020~\cite{isotopes-nubase2020}; binding control: AME2020~\cite{isotopes-ame2020}.}

\subsection{A state-dependent finite law}
\label{manual2:survival-law}
\paragraph{Setup.}
The third application model instantiates the conditional H1 kernel of
\mainequation{eq:admitted-kernel}, paired orientation response of
\mainequation{eq:sadar-orientation-antisymmetry}, and excess construction of
\mainequation{eq:shedding-excess}. Main does not select a nuclear state, its kernel weights
or a physical loss observation for this example. Within the declared finite scope, the
application fixes two retained square paths
\begin{align}
\gamma_+&=(0000,1000,1001,0001,0000),\\
\gamma_-&=(0000,0100,0101,0001,0000).
\label{eq:manual2-survival-paths}
\end{align}
Their vertices have one injective Q4 encoding. At \(0000\), the incoming ancestor slot
is \(0001\); the three continuations are \(1000,0100,0010\). The third exits through
\(0010\), then takes \(1010\) or \(0110\). These last routes carry hypothetical alpha
and fission labels, each with declared conditional weight one. Their nuclear identification
is not established.

A separate reference lineage repeats \(\gamma_+\). Both lineages begin with two distinct
retained-start witnesses. One step of the survival experiment completes a new reference
return, binds the already committed candidate prefix, normalizes the kernel, executes the
selected path and commits the response for the next contact. Four H1 traversals and the
retained start witness the finite square return; passing this predicate does not establish
native motif admission. The structural hinge is retained. No universal seconds per
traversal or scheduler cut are assigned.

\paragraph{Derivation.}
Let \(h\) be a particular retained choice history, with signed committed memory \(R_h\),
common capacity \(C\ge0\), and common positive loss gain \(\lambda\). The declared
specialization is
\begin{align}
X_h&=\max(0,|R_h|-C),&
(w_+,w_-,w_{\rm exit})&=(1,1,\lambda X_h),\\
P_+(h)=P_-(h)&=\frac{1}{2+\lambda X_h},&
P_{\rm exit}(h)&=\frac{\lambda X_h}{2+\lambda X_h},\label{eq:manual2-survival-kernel}\\
R_{h\sigma}&=R_h+\sigma,&\sigma&\in\{-1,+1\}.
\end{align}
The last update retains the oriented response pair \((\sigma,-\sigma)\). Its unit
magnitude specializes declared pressure \(2\) and signed lag magnitude \(1/2\);
neither is an isotope-derived support accessor or a lag inferred from measured decay.
This is the explicit application premise between the source response form and the
weight law. Zero exit weight leaves the exit structurally admitted and unrealized.

For exact history mass \(\mu_h\), multiplication by the probabilities in
\eqnref{eq:manual2-survival-kernel} supplies the child masses. Absorbed histories retain
their completed return IDs and inner scope map. They have no later candidate return.
Distinct prefixes remain distinct even if their numerical memories agree. Summing only
retained histories after reference return \(c\) gives
\begin{equation}
S(c)=\Pr(C_{\rm loss}>c)=\sum_{h\ \mathrm{retained\ after}\ c}\mu_h,
\qquad
c_{1/2}=\min\{c:S(c)\le\tfrac12\}.
\label{eq:manual2-survival-median}
\end{equation}
Thus \(c_{1/2}\) is the first discrete crossing. The recorded crossing interval
\((c_{1/2}-1,c_{1/2}]\) makes no within-cycle interpolation. If no crossing occurs by
the horizon, the result is a censored lower bound. A single loss event is not this
ensemble median.

\paragraph{Specialization.}
With \(R_\varnothing=0,C=1,\lambda=1\), exact propagation gives
\begin{equation}
\bigl(S(0),\ldots,S(8)\bigr)=
\left(1,1,1,\frac56,\frac34,\frac35,\frac{19}{36},
\frac{703}{1680},\frac{211}{576}\right).
\label{eq:manual2-survival-results}
\end{equation}
The median crosses in \((6,7]\) reference returns. Each run expands 255 retained
prefixes without merging. Doubling the common loss gain gives \((4,5]\); setting
the common capacity to zero gives \((2,3]\). Disabling new response commits at return
2 leaves \(S(c)=1\) through return 8. Its survival agrees with the baseline through
return 2 and first differs at return 3: \(1\) versus \(5/6\). This controlled change
therefore affects subsequent loss rather than merely annotating a precomputed trace.

The finite law is evaluated with the exact M3 mapping implementation and
\mainequation{eq:m3-finite-law}. For the retained indicator at return 8, its expectation
is \(211/576\) and its variance is \((211/576)(365/576)\). This ensemble spread is
distinct from uncertainty about the nuclear model or the experimental calibration.

\paragraph{Calibration.}
No isotope-specific hazard or capacity is fitted. Every named isotope currently labels
the same surrogate configuration; \eqnref{eq:manual2-isotope-grouping} does not change
its law. A physical elapsed-time report would require an observation map taking reference
returns to observed oscillator cycles. For example, the hypothetical constant gain
\(g>0\) would give \(T=gc/\Delta\nu_{\rm Cs}\), with
\(\Delta\nu_{\rm Cs}=9\,192\,631\,770\,\mathrm{Hz}\). The unselected gains \(g=1,2\)
give different times while preserving all native histories and the same SI definition.
The defining frequency alone therefore cannot select a decay-time prediction. Realization,
model discrepancy and covariance remain unresolved rather than numerical zeros.

\paragraph{Falsification.}
The exact tests reject probability leakage, coordinate aliases, lost lineage, a forged
return, nonantisymmetric response, later survival disconnected from feedback, and status
promotion. The target-free prediction runs without the isotope catalog. The comparison
replays that complete prediction before consuming validated source records. A changed
target or a self-consistent replacement hash cannot rewrite the frozen prediction.

\aodprovenance{\path{manual2/data/m1/survival/} records the common model contract, four
exact runs, prediction, comparison and unresolved construction obligations. The earlier
finite folding and stateful fixtures remain separate. The observation-gain dependency
instantiates \mainsection{prop:m3-atomic-gain}.}
\aodliteraturenote{SI reference definition: BIPM~\cite{bipm-second}.}

\subsection{A measured test of the restricted specialization}
\paragraph{Setup.}
Equal surrogate states with a common observation and clock map imply identical survival
distributions and a half-life ratio of one. This restricted consequence is testable even
though the isotope-to-native-state construction is incomplete. The selected adopted
records were exposed before this calculation; this is a retrospective comparison.

\paragraph{Derivation.}
Using numerator and denominator source uncertainty envelopes without treating them as
independent Gaussian inputs yields
\begin{center}
\begin{tabular}{lcc}
\toprule
Half-life ratio & Adopted central ratio & Source-envelope ratio range\\
\midrule
Og-294 / Lv-290 & \(29/415\) & \([2/59,51/320]\)\\
Lv-290 / Fl-286 & \(1/20\) & \([16/515,59/695]\)\\
Lv-293 / Fl-289 & \(19/480\) & \([17/800,79/950]\)\\
\bottomrule
\end{tabular}
\end{center}
All three ranges exclude one. These are descriptive enclosures constructed from the
reported uncertainties, not a joint confidence statement or significance score.

\paragraph{Specialization.}
The identity-insensitive, shared-clock specialization disagrees with these reported
envelopes. A successful common law must therefore distinguish the isotope configurations
or their justified observation relations. The measurements do not select a replacement
law. This result does not falsify the general AOD construction or establish a native
nuclear prediction. The next obligations are an admitted isotope/state configuration,
an independently justified state-to-loss law, the reference observation map and a complete
physical uncertainty model. Fl-298 remains an unscored target.

\paragraph{Calibration.}
A separate conventional benchmark imports each adopted \(T_{1/2}\) and defines
\(S(t)=\exp(-\lambda_{\rm exp}t)\),
\(\lambda_{\rm exp}=\ln2/T_{1/2}\). Its mean lifetime is
\(T_{1/2}/\ln2\), not \(T_{1/2}\). Reproducing survival \(1/2\) at an imported
half-life checks the conventional calculation; it does not predict that half-life.
The native finite law in \eqnref{eq:manual2-survival-kernel} is not assumed exponential.

\paragraph{Falsification.}
A physical isotope prediction must fix its state construction and observation law before
joining the evaluation target. Independent prediction additionally requires an exposure
design beyond these already known records. Source-specific revisions remain separately
bound comparisons, and no unobserved lifetime can be counted as agreement.

\aodprovenance{The separate \path{prediction.json} and \path{comparison.json} bind their
complete inputs and exact outputs. Tests SURV-01--19 state their finite assertion scope;
passing them does not certify nuclear admission.}
